\section{讨论}
\label{sec:discussion}

\textbf{从通用 VLA 到农业特需 VLA。}本文的结果表明，通用 VLA 在桌面任务上形成的记忆设计（固定帧数、按时间均匀采样）在农业长时作业中存在结构性缺陷，而且这一缺陷无法通过调参或适度扩大预算消除。农业场景真正需要的不是“记更多帧”，而是“记对的东西”：场景锚保住被遮挡之前的画面，工作事件层保住刚刚发生的细节，摘要层以极低的代价保住整行作业的进度。HarvestMem 只替换了通用结构中“记什么、怎么组织”这一环节，其余部分原样沿用，因此可以直接嫁接到不同的 VLA 骨干上，并随通用模型的进步而进步。我们认为，这种“通用骨干 + 领域特需记忆”的范式，是把大模型能力落地到农业装备的一条可行路径。

\textbf{对农业作业的意义。}按配额采收、分级抽样、疏果和按订单采摘，本质上都要求机器人记住自己做过什么；在枝叶茂密的冠层中，“曾经看见过什么”同样需要被记住。HarvestMem 不改变基础模型的记忆帧数和推理开销，新增参数仅约 \V{hem_params}~M，\pedet{} 是可与控制循环并行的轻量网络，事后确认自训练只依赖夹爪和腕部相机这两种所有采摘机器人都具备的传感器。对算力受限、难以人工标注的田间平台，这比扩大记忆预算、引入大型记忆模块或在植株上布设传感器都更容易落地。

\textbf{为什么“无差异”也是重要结果。}遮挡 v2 上各记忆条件等效，并不说明方法失效，而是确认了它的作用边界：场景信息集中在开头画面中，均匀采样与场景锚层都会保留它；\hem{} 只改变事件型信息的保留方式，并且没有挤掉开头画面。正因为两类任务的预测在评测前就已写定，事件记忆任务上的增益才更可信，而不是事后挑出的有利结果。同理，计数 v2 上 HEM 与 K8 的差距较小、FruitRow 上差距显著，也都与容量分析的事先预测一致。

\textbf{基准设计的教训。}v1 的经验表明，按直觉设计的“记忆任务”很可能存在捷径：环境的终止逻辑、遮挡物的物理交互、甚至成熟度的先验比例，都可能让策略绕开记忆。我们建议，任何声称检验记忆的机器人基准都应报告信息受限脚本策略的审计结果，并设置正对照。这一做法成本很低，却能决定实验结论是否有意义。

\textbf{为什么需要学习型检测器与事后确认。}几何距离只描述“靠近了果实”，而采摘任务关心的是“果实是否真的被摘下”；手工感知链路证明了“从感知中恢复事件”是可行的，也暴露了它对颜色提示和光照的依赖。\pedet{} 把检测目标从几何量换成物理结果，使闭环中常见的失败尝试不再污染记忆。事后确认自训练进一步解决了真机标签的来源问题：采摘的成败在果实入篮时才可确认，确认后回溯标注事件时刻，既不需要人工标注，也不需要在果实上安装任何装置，检测器因此可以随作物、季节和光照的变化持续更新。

\textbf{局限。}第一，仿真中的果实为刚体球，果柄约束经过简化，真实果实的形变、果柄柔性和成簇生长尚未建模。第二，真机主实验采用“真实番茄 + 可复位植株骨架”，以换取可重复的配对比较，与温室和果园中的枝叶柔性、自然风扰和光照变化仍有差距；温室验证规模较小，不做配对检验；用窄带照明模拟“颜色不可分”也只是场景记忆需求的一种实验室近似。第三，事件定义针对“夹持—拉拽”式采摘，剪切式末端执行器需要把“剪刀闭合”纳入事件定义；事后确认依赖“入篮”这一结果，若采用传送带或吸管式收集，需要相应修改确认规则。第四，拼图帧以较低分辨率保存远期事件，若任务需要回忆远期事件的细节（如某颗果的确切成熟度），摘要层可能不够用。第五，真机每条件 \V{rb_trials} 次试验只能检测较大的效应，较小的差异需要更大规模的田间验证。

\textbf{未来工作。}下一步计划在温室番茄和果园苹果上开展移动平台的整行连续作业试验，将事件语义扩展到放果确认、掉果检测、剪梗与抓取重试，并研究把事件记忆与三维场景记忆结合，使机器人同时记住“哪里有什么”和“自己做过什么”。

\section{结论}
\label{sec:conclusion}
本文面向农业长时采摘作业，研究了 VLA 模型在固定记忆预算下“记什么、怎么组织”的问题。理论分析和实测表明，通用 VLA 按时间均匀采样的记忆对短暂采摘事件的覆盖率随作业时长按 $1/t$ 衰减，16 帧记忆平均只记得约 \V{cap_u} 次采摘。为此提出事件中心分层记忆 VLA 模型 HarvestMem，包含三项创新：物理监督的学习型采摘事件检测器 \pedet，零人工标注训练，能区分成功采摘与失败尝试并对光照鲁棒；事件中心分层记忆 \hem，以场景锚、工作事件、情景摘要拼图帧和时间上下文组织固定预算，并以类型-时间嵌入编码记忆语义，在不增加推理开销的前提下把可保留的采摘次数提升到 \V{cap_hem} 次；事件扰动训练与事后确认自训练，使策略对检测误差鲁棒，并使检测器无需人工标注、无需附加传感器即可适配真机。在经过捷径审计的 \bench{} 基准上，HarvestMem 把事件记忆计数任务成功率从 \pct{c_k0_s} 提高到 \pct{c_hem_s}，把多株长时配额任务从 \pct{l_k0_s} 提高到 \pct{l_hem_s}，在场景记忆任务上与均匀记忆等效，与事先预测一致；真实番茄计数任务成功率从 \pct{rc_k0_s} 提高到 \pct{rc_hem_s}。按物理事件组织记忆，为通用 VLA 走向农业长时作业提供了一条简单、可审计且易于部署的途径。
